% 사과 가격 예측 선형 회귀 시뮬레이터 이론해설서 v1.4.0.1
% LuaLaTeX 컴파일 필요 (컬러 이모지 + 한글 지원)

\documentclass[11pt,a4paper]{article}

% === 패키지 ===
\usepackage{fontspec}
\usepackage{kotex}
\usepackage{geometry}
\usepackage{graphicx}
\usepackage{xcolor}
\usepackage{tcolorbox}
\usepackage{booktabs}
\usepackage{array}
\usepackage{longtable}
\usepackage{colortbl}
\usepackage{fancyhdr}
\usepackage{titlesec}
\usepackage{hyperref}
\usepackage{emoji}

% === 페이지 설정 ===
\geometry{left=25mm, right=25mm, top=30mm, bottom=30mm}

% === 폰트 설정 ===
\setmainfont{Noto Sans CJK KR}
\setsansfont{Noto Sans CJK KR}
\setmonofont{Noto Sans CJK KR}
\setemojifont{Noto Color Emoji}

% === 색상 정의 ===
\definecolor{primary}{HTML}{E53935}
\definecolor{darkbrown}{HTML}{3E2723}
\definecolor{mediumbrown}{HTML}{5D4037}
\definecolor{lightbrown}{HTML}{8D6E63}
\definecolor{cream}{HTML}{FFF3E0}
\definecolor{lightgray}{HTML}{F5F5F5}
\definecolor{tableheader}{HTML}{E8E8E8}

% === 제목 스타일 ===
\titleformat{\section}
  {\Large\bfseries\color{primary}}
  {\thesection.}{0.5em}{}
\titleformat{\subsection}
  {\large\bfseries\color{darkbrown}}
  {\thesubsection}{0.5em}{}

% === 머리글/바닥글 ===
\pagestyle{fancy}
\fancyhf{}
\fancyhead[R]{\small\color{lightbrown}선형 회귀 이론해설서 v1.4.0.1}
\fancyfoot[C]{\thepage}
\renewcommand{\headrulewidth}{0.4pt}

% === tcolorbox 스타일 ===
\tcbset{
  tipbox/.style={
    colback=cream,
    colframe=primary,
    boxrule=1.5pt,
    arc=3mm,
    left=5mm,
    right=5mm,
    top=3mm,
    bottom=3mm,
    fonttitle=\bfseries\color{primary}
  },
  formulabox/.style={
    colback=lightgray,
    colframe=lightgray,
    boxrule=0pt,
    arc=2mm,
    left=5mm,
    right=5mm,
    top=3mm,
    bottom=3mm
  },
  summarybox/.style={
    colback=white,
    colframe=primary,
    boxrule=2pt,
    arc=4mm,
    left=5mm,
    right=5mm,
    top=5mm,
    bottom=5mm
  }
}

% === 수식 해석 환경 ===
\newcommand{\formulawithexp}[2]{%
  \begin{tcolorbox}[formulabox]
    \centering\large\texttt{#1}\\[3mm]
    \small\color{mediumbrown} \emoji{point-right} #2
  \end{tcolorbox}
}

% === 하이퍼링크 설정 ===
\hypersetup{
  colorlinks=true,
  linkcolor=primary,
  urlcolor=primary
}

\begin{document}

% 목차 제목 한글화
\renewcommand{\contentsname}{목차}

% ===== 표지 =====
\thispagestyle{empty}
\begin{center}
\vspace*{3cm}

{\fontsize{72pt}{72pt}\selectfont \emoji{red-apple}}

\vspace{1cm}

{\fontsize{28pt}{32pt}\selectfont\bfseries\color{primary} 선형 회귀 이론해설서}

\vspace{0.5cm}

{\Large\color{lightbrown} Linear Regression Theory Guide}

\vspace{1.5cm}

\begin{tcolorbox}[colback=primary, colframe=primary, arc=15pt, width=4cm]
  \centering\large\bfseries\color{white} v1.4.0.1
\end{tcolorbox}

\vspace{1cm}

{\large 품질 점수 기반 단순 선형 회귀}

{\color{lightbrown} 시뮬레이터 v1.2.4.1 연동}

\vspace{3cm}

{\color{lightbrown} 청주교육대학교 컴퓨팅사고 교육자료}

{\color{lightbrown} 2026년 2월}

\end{center}

\newpage

% ===== 목차 =====
\tableofcontents
\newpage

% ===== 30초 퀵 서머리 (신규) =====
\section*{\emoji{stopwatch} 30초 퀵 서머리}
\addcontentsline{toc}{section}{30초 퀵 서머리}

\begin{tcolorbox}[summarybox, title={\large\emoji{light-bulb} 이 문서의 핵심을 한눈에!}]

\begin{center}
\renewcommand{\arraystretch}{1.6}
\begin{tabular}{|>{\centering\arraybackslash}p{4cm}|>{\centering\arraybackslash}p{4cm}|>{\centering\arraybackslash}p{4cm}|}
\hline
\rowcolor{tableheader}
\textbf{우리가 하는 일} & \textbf{AI 용어로 바꾸면?} & \textbf{비유하자면?} \\
\hline
품질 점수 만들기 & \textbf{특성 추출 / 정규화} & \emoji{kitchen-knife} 재료 손질하기 \\
\hline
가격 예측 공식 찾기 & \textbf{모델 학습} & \emoji{books} 정답 맞히기 연습 \\
\hline
얼마나 틀렸나 계산 & \textbf{손실 함수 (MSE)} & \emoji{cross-mark} 오답 체크하기 \\
\hline
정답에 가깝게 수정 & \textbf{경사 하강법} & \emoji{mountain} 눈감고 골짜기 찾기 \\
\hline
잘 맞추는지 점수 매김 & \textbf{R² 평가} & \emoji{scroll} AI 성적표 \\
\hline
새 사과 가격 맞추기 & \textbf{추론 (Inference)} & \emoji{crystal-ball} 시험 보기 \\
\hline
\end{tabular}
\end{center}

\vspace{5mm}

\begin{center}
\textbf{\large \emoji{red-apple} 핵심 공식}
\end{center}

\begin{tcolorbox}[formulabox]
\centering\Large
\texttt{가격 = w × 품질 + b}
\end{tcolorbox}

\begin{itemize}
  \item \textbf{w (가중치)}: 품질이 1점 오를 때 가격이 얼마나 오르는가
  \item \textbf{b (편향)}: 품질이 0점일 때의 기본 가격
\end{itemize}

\end{tcolorbox}

\newpage

% ===== 1. 머신러닝이란? =====
\section{머신러닝이란?}

머신러닝(Machine Learning)은 컴퓨터가 명시적으로 프로그래밍되지 않고도 \textbf{데이터로부터 학습}하여 예측이나 결정을 내릴 수 있게 하는 인공지능의 한 분야입니다.

\subsection{머신러닝의 종류}

\begin{center}
\renewcommand{\arraystretch}{1.4}
\begin{tabular}{|c|p{6cm}|c|}
\hline
\rowcolor{tableheader}
\textbf{유형} & \textbf{설명} & \textbf{시뮬레이터 예시} \\
\hline
지도 학습 & 입력과 정답(레이블)이 함께 제공됨 & 품질→가격 예측 \emoji{check-mark-button} \\
\hline
비지도 학습 & 정답 없이 데이터의 패턴을 찾음 & 군집화 (K-means) \\
\hline
강화 학습 & 보상을 최대화하는 행동을 학습 & 게임 AI (알파고) \\
\hline
\end{tabular}
\end{center}

\begin{figure}[h]
\centering
\includegraphics[width=0.8\textwidth]{theory_images/img01_ml_types.png}
\caption{머신러닝의 종류}
\end{figure}

\begin{tcolorbox}[tipbox, title={\emoji{light-bulb} 시뮬레이터에서 확인하기}]
사과 가격 예측 시뮬레이터는 \textbf{'지도 학습'} 중 \textbf{'회귀'} 문제를 다룹니다. 품질(입력)과 가격(정답)이 함께 제공된 데이터로 학습합니다.
\end{tcolorbox}

\newpage

% ===== 2. 선형 회귀 =====
\section{선형 회귀 (Linear Regression)}

선형 회귀는 입력 변수(특성)와 출력 변수(타겟) 사이의 \textbf{선형 관계}를 모델링하는 지도 학습 알고리즘입니다.

\subsection{단순 선형 회귀 수식}

\formulawithexp{y = w × x + b}{예측 가격 = (품질의 영향력 × 품질 점수) + 기본 가격}

\begin{center}
\renewcommand{\arraystretch}{1.4}
\begin{tabular}{|c|c|p{5.5cm}|}
\hline
\rowcolor{tableheader}
\textbf{기호} & \textbf{의미} & \textbf{시뮬레이터 예시} \\
\hline
y & 예측값 (출력) & 예측 가격 (원) \\
\hline
x & 입력값 (특성) & 품질 점수 (0\textasciitilde100) \\
\hline
w & 가중치 (기울기) & 품질이 가격에 미치는 영향력 \\
\hline
b & 편향 (y절편) & 품질이 0일 때 기본 가격 \\
\hline
\end{tabular}
\end{center}

\begin{figure}[h]
\centering
\includegraphics[width=0.7\textwidth]{theory_images/img03_linear.png}
\caption{단순 선형 회귀 모델}
\end{figure}

\subsection{시뮬레이터의 회귀 모델}

\begin{tcolorbox}[formulabox]
\centering\Large
\texttt{가격 = w × 품질 + b}
\end{tcolorbox}

시뮬레이터는 당도와 무게를 \textbf{품질 점수}로 통합한 후, 1변수 단순 선형 회귀를 적용합니다.

\newpage

% ===== 3. 품질 점수 계산 =====
\section{품질 점수 계산 (특성 통합)}

여러 특성을 하나로 통합하면 시각화가 단순해지고 학습 개념을 이해하기 쉬워집니다.

\begin{tcolorbox}[formulabox]
\centering\Large
\texttt{품질 = 당도점수(50\%) + 무게점수(50\%)}
\end{tcolorbox}

\subsection{세부 계산식}

\begin{center}
\renewcommand{\arraystretch}{1.4}
\begin{tabular}{|c|c|c|}
\hline
\rowcolor{tableheader}
\textbf{구성요소} & \textbf{계산식} & \textbf{범위} \\
\hline
당도점수 & (당도 - 10) / 8 × 50 & 10\textasciitilde18 Brix → 0\textasciitilde50점 \\
\hline
무게점수 & (무게 - 150) / 200 × 50 & 150\textasciitilde350g → 0\textasciitilde50점 \\
\hline
품질 (총합) & 당도점수 + 무게점수 & 0\textasciitilde100점 \\
\hline
\end{tabular}
\end{center}

\subsection{계산 예시}

당도 14 Brix, 무게 250g인 사과:

\begin{itemize}
  \item 당도점수 = (14 - 10) / 8 × 50 = \textbf{25점}
  \item 무게점수 = (250 - 150) / 200 × 50 = \textbf{25점}
  \item 총 품질 = 25 + 25 = \textbf{50점}
\end{itemize}

\begin{figure}[h]
\centering
\includegraphics[width=0.7\textwidth]{theory_images/img02_quality.png}
\caption{품질 점수 계산 과정}
\end{figure}

\newpage

% ===== 4. 정규화 =====
\section{정규화 (Normalization)}

정규화는 서로 다른 범위를 가진 데이터를 일정한 범위로 변환하는 과정입니다.

\subsection{왜 정규화가 필요한가?}

\begin{tcolorbox}[tipbox, title={\emoji{straight-ruler} 비유: 서로 다른 자(Ruler) 합치기}]
당도는 10\textasciitilde18 Brix, 무게는 150\textasciitilde350g입니다.

만약 그냥 더하면? 무게(최대 350)가 당도(최대 18)보다 \textbf{20배나 커서}, 무게만 가격에 영향을 주는 것처럼 보입니다.

이를 공평하게 0\textasciitilde100 사이로 맞추는 과정이 정규화입니다. 마치 \textbf{cm 자와 km 자를 모두 \% 자로 바꾸는 것}과 같습니다!
\end{tcolorbox}

\begin{figure}[h]
\centering
\includegraphics[width=0.85\textwidth]{theory_images/img09_normalization_analogy.png}
\caption{정규화 비유: 서로 다른 단위의 자를 통일된 \% 자로 변환}
\end{figure}

\begin{center}
\renewcommand{\arraystretch}{1.4}
\begin{tabular}{|p{5.5cm}|p{5.5cm}|}
\hline
\rowcolor{tableheader}
\textbf{\emoji{cross-mark} 정규화 없이} & \textbf{\emoji{check-mark-button} 정규화 적용} \\
\hline
학습이 불안정해질 수 있음 & 안정적이고 빠른 학습 가능 \\
\hline
가중치가 극단적인 값으로 발산 & 가중치가 적절한 범위 유지 \\
\hline
학습률 설정이 어려움 & 학습률 선택이 용이 \\
\hline
\end{tabular}
\end{center}

\subsection{Z-score 정규화 수식}

\formulawithexp{정규화된 값 = (원래 값 - 평균) / 표준편차}{``평균에서 얼마나 떨어졌는가''를 표준 단위로 표현}

정규화된 데이터는 \textbf{평균이 0, 표준편차가 1}인 분포를 가집니다.

\begin{tcolorbox}[tipbox, title={\emoji{magnifying-glass-tilted-left} 시뮬레이터 내부 동작}]
시뮬레이터는 정규화된 가중치(w, b)로 학습하고, 결과 표시 시 실제 스케일의 가중치(wReal, bReal)로 변환하여 보여줍니다.
\end{tcolorbox}

\newpage

% ===== 5. 손실 함수 =====
\section{손실 함수 (Loss Function)}

\textbf{손실(Loss)}은 모델의 예측값과 실제값의 차이를 수치화한 것입니다.

\subsection{평균 제곱 오차 (MSE)}

\formulawithexp{MSE = (1/n) × Σ(실제값 - 예측값)²}{``오차들을 제곱해서 평균 낸 값'' - 큰 오차에 더 큰 페널티!}

MSE는 오차를 제곱하여 평균을 구합니다. \textbf{큰 오차에 더 큰 페널티}를 부여합니다.

\subsection{오차(잔차)의 의미}

\formulawithexp{오차 = 실제 가격 - 예측 가격}{``AI가 틀린 정도'' - 양수면 과소예측, 음수면 과대예측}

\subsection{오차선 시각화}

'오차선 표시' 기능은 각 데이터 포인트에서 회귀선까지의 거리를 시각화합니다:

\begin{itemize}
  \item \textbf{점선}: 실제 데이터 점 ↔ 회귀선 위 예측점
  \item \textbf{색상 변화}: \textcolor{green!70!black}{녹색(작은 오차)} → \textcolor{red}{빨강(큰 오차)}
\end{itemize}

\begin{figure}[h]
\centering
\includegraphics[width=0.7\textwidth]{theory_images/img04_loss.png}
\caption{손실 함수와 오차선 시각화}
\end{figure}

\newpage

% ===== 6. 경사 하강법 =====
\section{경사 하강법 (Gradient Descent)}

경사 하강법은 손실 함수의 기울기를 따라 손실이 감소하는 방향으로 가중치를 업데이트합니다.

\begin{tcolorbox}[tipbox, title={\emoji{mountain} 비유: 안개 낀 산에서 골짜기 찾기}]
안개가 너무 짙어 \textbf{발 앞만 보이는 산}에서 가장 낮은 골짜기(오차가 가장 작은 지점)를 찾는 과정입니다.

한 발 내딛어보고(\textbf{학습률}), 경사가 아래쪽이면 그 방향으로 계속 내려갑니다.

\begin{itemize}
  \item \emoji{warning} 너무 크게 걸으면 → 골짜기를 지나쳐버림 (발산)
  \item \emoji{snail} 너무 작게 걸으면 → 내려오는 데 하루 종일 (느린 수렴)
\end{itemize}
\end{tcolorbox}

\begin{figure}[h]
\centering
\includegraphics[width=0.85\textwidth]{theory_images/img10_gradient_analogy.png}
\caption{경사 하강법 비유: 안개 낀 산에서 발밑만 보고 골짜기 찾기}
\end{figure}

\subsection{업데이트 공식}

\formulawithexp{w\_new = w\_old - 학습률 × 경사방향}{``새 가중치 = 기존 위치 - (한 걸음 크기 × 내려갈 방향)''}

\begin{figure}[h]
\centering
\includegraphics[width=0.65\textwidth]{theory_images/img05_gradient.png}
\caption{경사 하강법 원리}
\end{figure}

\subsection{학습률의 영향}

\begin{center}
\renewcommand{\arraystretch}{1.4}
\begin{tabular}{|c|c|c|}
\hline
\rowcolor{tableheader}
\textbf{학습률} & \textbf{특징} & \textbf{손실 그래프} \\
\hline
너무 작음 (0.001) & 수렴 매우 느림 & 완만하게 천천히 감소 \\
\hline
\textbf{적절함 (0.01)} \emoji{check-mark-button} & 안정적 수렴 & 빠르고 안정적으로 감소 \\
\hline
너무 큼 (0.1) & 발산 가능 & 급격히 증가하거나 진동 \\
\hline
\end{tabular}
\end{center}

\begin{figure}[h]
\centering
\includegraphics[width=0.8\textwidth]{theory_images/img06_lr.png}
\caption{학습률에 따른 손실 그래프 비교}
\end{figure}

\begin{tcolorbox}[tipbox, title={\emoji{video-game} 시뮬레이터에서 확인하기}]
\emoji{play-button} 재생 버튼을 클릭하면 회귀선이 데이터에 맞춰가는 애니메이션을 볼 수 있습니다!
\end{tcolorbox}

\newpage

% ===== 7. 모델 평가 =====
\section{모델 평가}

\subsection{결정 계수 (R²)}

\formulawithexp{R² = 1 - (잔차 제곱합 / 총 제곱합)}{``모델이 데이터의 변동성을 얼마나 잘 설명하는가'' (1에 가까울수록 좋음)}

\begin{center}
\renewcommand{\arraystretch}{1.4}
\begin{tabular}{|c|c|c|}
\hline
\rowcolor{tableheader}
\textbf{R² 값} & \textbf{의미} & \textbf{시뮬레이터 표시} \\
\hline
R² = 1 & 완벽한 적합 & - \\
\hline
R² > 0.8 & 모델이 데이터를 잘 설명함 & \emoji{check-mark-button} 양호 \\
\hline
R² > 0.5 & 보통 수준의 설명력 & \emoji{warning} 보통 \\
\hline
R² < 0 & 평균보다 못한 예측 & \emoji{cross-mark} 부족 \\
\hline
\end{tabular}
\end{center}

\begin{figure}[h]
\centering
\includegraphics[width=0.65\textwidth]{theory_images/img07_r2.png}
\caption{R² 점수 해석 가이드}
\end{figure}

\begin{tcolorbox}[tipbox, title={\emoji{scroll} R²는 AI의 성적표!}]
학교 시험에서 100점 만점에 80점을 받으면 ``잘했다''고 하듯이, R² > 0.8이면 AI가 ``공부를 잘 했다''는 뜻입니다!
\end{tcolorbox}

\newpage

% ===== 8. 추론 =====
\section{추론 (Inference)}

추론은 학습된 모델을 사용하여 \textbf{새로운 데이터에 대한 예측}을 수행하는 과정입니다.

\subsection{추론 과정 4단계}

\begin{enumerate}
  \item \emoji{honey-pot} \emoji{balance-scale} 당도, 무게 슬라이더로 새 사과 특성 입력
  \item \emoji{star} 품질 점수 자동 계산 (당도점수 + 무게점수)
  \item \emoji{abacus} 학습된 회귀식에 품질 점수 대입
  \item \emoji{money-bag} 예측 가격 출력 및 차트에 시각화
\end{enumerate}

\subsection{예측 시각화}

추론 모드에서는 다음 요소들이 차트에 표시됩니다:

\begin{itemize}
  \item \textbf{수직 점선}: 품질 점수 위치에서 회귀선까지
  \item \textbf{수평 점선}: 회귀선에서 가격 축까지
  \item \textbf{노란 원} \emoji{yellow-circle}: 예측 위치 (품질, 가격)
\end{itemize}

\begin{tcolorbox}[tipbox, title={\emoji{crystal-ball} 자동 추론 기능}]
슬라이더를 움직이면 \textbf{즉시} 예측 결과가 업데이트됩니다. 실시간으로 ``만약 이런 사과라면?''을 실험해 볼 수 있습니다!
\end{tcolorbox}

\newpage

% ===== 9. 학습 과정 요약 =====
\section{학습 과정 요약}

\begin{center}
\renewcommand{\arraystretch}{1.4}
\begin{tabular}{|c|c|p{6cm}|}
\hline
\rowcolor{tableheader}
\textbf{단계} & \textbf{과정} & \textbf{설명} \\
\hline
1 & 데이터 준비 & 당도, 무게 → 품질 점수 계산 \\
\hline
2 & 가중치 초기화 & w=0, b=0으로 시작 \\
\hline
3 & 예측 & 가격 = w × 품질 + b 계산 \\
\hline
4 & 손실 계산 & MSE 계산 (오차선 시각화) \\
\hline
5 & 경사 방향 확인 & 손실이 줄어드는 방향 파악 \\
\hline
6 & 가중치 업데이트 & w, b 조정 (회귀선 이동) \\
\hline
7 & 반복 & 에포크 수만큼 3\textasciitilde6단계 반복 \\
\hline
\end{tabular}
\end{center}

\begin{figure}[h]
\centering
\includegraphics[width=0.85\textwidth]{theory_images/img08_flow.png}
\caption{선형 회귀 학습 과정 흐름도}
\end{figure}

\newpage

% ===== 10. 교육적 포인트 =====
\section{교육적 포인트: AI는 항상 맞지 않는다}

'오차선 표시' 기능은 AI의 한계를 시각적으로 보여주는 중요한 교육 도구입니다.

\subsection{오차선이 보여주는 것}

\begin{itemize}
  \item 모든 예측에는 \textbf{오차가 존재}함
  \item 같은 품질이라도 실제 가격은 다를 수 있음
  \item 노이즈(현실 세계의 불확실성)를 완벽히 예측할 수 없음
\end{itemize}

\subsection{수업에서의 활용}

\begin{enumerate}
  \item 노이즈 수준을 5\% → 20\%로 변경하며 R² 변화 관찰
  \item 오차선 색상(\textcolor{green!70!black}{녹색}→\textcolor{red}{빨강})으로 오차 크기 직관적 이해
  \item ``완벽한 예측은 불가능하다''는 AI의 한계 토론
\end{enumerate}

\subsection{데이터가 AI의 실력을 결정한다 (신규)}

\begin{tcolorbox}[tipbox, title={\emoji{red-apple} 생각해보기: 빨간 사과 vs 초록 사과}]
만약 우리가 \textbf{``빨간 사과''} 데이터만 가지고 학습했다면, \textbf{``초록 사과(아오리)''의 가격}도 잘 맞출 수 있을까요?

\vspace{3mm}
\emoji{point-right} AI의 실력은 알고리즘만큼이나 \textbf{``어떤 데이터를 주었는가''}에 달려 있습니다!
\end{tcolorbox}

\begin{figure}[h]
\centering
\includegraphics[width=0.9\textwidth]{theory_images/img11_data_bias.png}
\caption{데이터 편향성: 빨간 사과로만 학습한 AI가 초록 사과를 예측할 때}
\end{figure}

\subsection{토론 질문}

\begin{tcolorbox}[colback=white, colframe=mediumbrown, boxrule=1pt, arc=2mm]
\emoji{books} \textbf{수업에서 활용할 토론 질문들}

\begin{enumerate}
  \item AI가 틀릴 수 있다면, 어떤 상황에서 AI를 신뢰해야 할까요?
  \item 오차를 줄이려면 어떤 데이터가 더 필요할까요?
  \item 학습 데이터에 없는 종류의 사과가 들어오면 어떻게 될까요?
  \item 계절에 따라 가격이 달라진다면 어떤 데이터가 필요할까요?
\end{enumerate}
\end{tcolorbox}

\newpage

% ===== 참고자료 =====
\section*{참고자료}
\addcontentsline{toc}{section}{참고자료}

\begin{center}
\renewcommand{\arraystretch}{1.3}
\begin{tabular}{|c|p{7.5cm}|c|c|}
\hline
\rowcolor{tableheader}
\textbf{\#} & \textbf{자료명} & \textbf{유형} & \textbf{비고} \\
\hline
1 & 사과가격예측\_선형회귀\_시뮬레이터\_v1.2.4.1.html & 파일 & 실습용 \\
\hline
2 & 사과가격예측\_선형회귀\_시뮬레이터\_개발일지\_v1.2.4.1.md & 문서 & 기술 상세 \\
\hline
3 & 사과가격예측\_선형회귀\_시뮬레이터\_퀵가이드\_v1.2.4.1.md & 문서 & 사용 안내 \\
\hline
4 & 이론해설서\_v1.2.0.0 (이전 버전) & 문서 & 수정 전 \\
\hline
5 & AI 실습 활동지 검토 의견 & 피드백 & 비유 추가 제안 \\
\hline
\end{tabular}
\end{center}

\vspace{1cm}

% ===== 생성/수정 이력 =====
\section*{생성/수정 이력}
\addcontentsline{toc}{section}{생성/수정 이력}

\begin{center}
\renewcommand{\arraystretch}{1.3}
\begin{tabular}{|c|c|c|c|p{4.5cm}|c|}
\hline
\rowcolor{tableheader}
\textbf{버전} & \textbf{날짜} & \textbf{시간} & \textbf{변경 수준} & \textbf{변경 내용} & \textbf{작성자} \\
\hline
v1.0.0.0 & 2026-01-31 & 14:00 & 최초 & 초안 작성 & Claude \\
\hline
v1.1.0.0 & 2026-01-31 & 15:40 & 마이너 & 품질점수, 단순선형회귀 반영 & Claude \\
\hline
v1.2.0.0 & 2026-02-02 & 14:00 & 마이너 & 정규화/추론/교육포인트 추가 & Claude \\
\hline
\textbf{v1.3.0.0} & \textbf{2026-02-02} & \textbf{15:30} & \textbf{마이너} & \textbf{퀵서머리, 비유, 수식해석 추가} & \textbf{Claude} \\
\hline
\textbf{v1.4.0.0} & \textbf{2026-02-02} & \textbf{14:35} & \textbf{리비전} & \textbf{비유 이미지 3개 삽입 (정규화, 경사하강법, 데이터편향)} & \textbf{Claude} \\
\hline
\textbf{v1.4.0.1} & \textbf{2026-02-02} & \textbf{14:45} & \textbf{패치} & \textbf{목차 제목 한글화 (Contents → 목차)} & \textbf{Claude} \\
\hline
\end{tabular}
\end{center}

\vspace{2cm}

\begin{center}
\textit{\color{lightbrown} 선형 회귀 이론해설서}
\end{center}

\end{document}
